\section*{\texorpdfstring{\S\CurrentExerciseGroup}{Section \CurrentExerciseGroup}}
\phantomsection
\addcontentsline{toc}{section}{\texorpdfstring{Exercises for \S\CurrentExerciseGroup}{Exercises for Section \CurrentExerciseGroup}}

\begin{enumerate}
\def\labelenumi{\arabic{enumi})}
\tightlist
\item
  \begin{enumerate}
  \def\labelenumii{\alph{enumii})}
  \tightlist
  \item
    Let \(\mathbf{K}\) be a compact convex subset of a Hausdorff locally convex space \(\mathbf{E}\) , \(\mu\) a positive measure on \(\mathbf{K}\) of mass 1, \(\mathbf{b}_{\mu}\) its barycenter. Show that for every positive function \(f\) on \(\mathbf{K}\) that is concave and lower semi-continuous, one has \(f(\mathbf{b}_{\mu}) \geqslant \int^{*} f d\mu\) . (Note that \(f\) is bounded (TVS, II, §2, Exer. 32), consequently \(-f\) is \(\mu\) -integrable.) From this, conclude that if \(f\) is a lower semi-continuous (or upper semi-continuous) function on \(\mathbf{K}\) that is both convex and concave, then \(\int f d\mu = f(\mathbf{b}_{\mu})\) .
  \end{enumerate}
\end{enumerate}

\begin{enumerate}
\def\labelenumi{\alph{enumi})}
\setcounter{enumi}{1}
\tightlist
\item
  Let I be the interval \([0,1]\) of R, K the subset of \(\mathcal{M}_{+}(I)\) formed by the positive measures of total mass 1, which is convex and compact for the vague topology, and let \(j: x \mapsto \varepsilon_{x}\) be the canonical injection of I into K, which is a homeomorphism of I onto a subspace of K. For every measure \(\nu \in K\) , set \(g(\nu) = \sum_{x \in I} \nu(\{x\})\) ; this is a function that
\end{enumerate}

is both convex and concave in K; moreover, if \(g_{n}(\nu)\) is the supremum of the \(\nu(A)\) for all the finite subsets A of I having at most n elements, then \(g_{n}\) is upper semi-continuous on K, and \(g(\nu) = \lim_{n \to \infty} g_{n}(\nu)\) for every \(\nu \in K\) , therefore g is \(\lambda\) -integrable for every measure \(\lambda\) on K. Let \(\mu\) be the measure on K such that \(\int f d\mu = \int_{\mathrm{I}} f(j(x)) dx\) for \(f \in \mathcal{K}(\mathrm{X}; \mathbf{R})\) , which is positive and of total mass 1; show that \(\mathbf{b}_{\mu}\) is Lebesgue measure on I and that \(\int g d\mu \neq g(\mathbf{b}_{\mu})\) .

\begin{enumerate}
\def\labelenumi{\arabic{enumi})}
\setcounter{enumi}{1}
\tightlist
\item
  Let X be a compact space, P a set of lower semi-continuous numerical functions on X, taking values in \(]-\infty,+\infty]\) . Assume that P contains the finite constants; for every function \(h\in P\) and every positive measure \(\mu\) on X, \(\mu^{*}(h)\) is defined and \(>-\infty\) (§4, Exer. 5).
\end{enumerate}

\begin{enumerate}
\def\labelenumi{\alph{enumi})}
\item
  For two positive measures \(\mu, \nu\) on \(X\) , write \(\mu \prec \nu\) if \(\mu^{*}(h) \leqslant \nu^{*}(h)\) for every function \(h \in \mathcal{P}\) . Show that the relation \(\mu \prec \nu\) is a pre-order relation on \(\mathcal{M}_{+}(X)\) ; it implies that \(\mu(1) = \nu(1)\) , \(c\mu \prec c\nu\) for all \(c > 0\) , and \(\mu + \lambda \prec \nu + \lambda\) for every measure \(\lambda \in \mathcal{M}_{+}(X)\) .
\item
  Assume that \(\mathcal{P} + \mathcal{P} \subset \mathcal{P}\) and that \(c \cdot \mathcal{P} \subset \mathcal{P}\) for \(c > 0\) . Let \(\mu\) be a positive measure on \(X\) , \(f\) a function in \(\mathcal{C}(X; \mathbf{R})\) . Denote by \(Q_f\) the set of functions \(h \geqslant f\) belonging to \(\mathcal{P}\) , and by \(M_\mu\) the set of measures \(\lambda \in \mathcal{M}_+(X)\) such that \(\lambda \prec \mu\) . Show that
\end{enumerate}

\[
\sup _ {\lambda \in \mathrm{M} _ {\mu}} \lambda (f) = \inf _ {h \in \mathrm{Q} _ {f}} \mu^ {*} (h).
\]

(For every function \(g \in \mathcal{C}(\mathrm{X};\mathbf{R})\) , let \(p(g) = \inf_{h \in \mathrm{Q}_g} \mu^*(h)\) . Show that \(p(g + g') \leqslant p(g) + p(g')\) and \(p(cg) = c \cdot p(g)\) for \(g, g'\) in \(\mathcal{C}(\mathrm{X};\mathbf{R})\) and \(c > 0\) ; prove on the other hand that \(\mathrm{M}_{\mu}\) is identical to the set of measures \(\lambda \in \mathcal{M}_{+}(\mathrm{X})\) such that, for every function \(g \in \mathcal{C}(\mathrm{X};\mathbf{R})\) , one has \(\lambda(g) \leqslant p(g)\) , and conclude by applying the Hahn-Banach theorem.) If the functions in \(\mathcal{P}\) are continuous and if \(S_0(f)\) is the lower envelope of \(Q_f\) , then also \(\mu(S_0(f)) = \sup \lambda(f)\) .

\[
\lambda \in \mathrm{M} _ {\mu}
\]

\begin{enumerate}
\def\labelenumi{\alph{enumi})}
\setcounter{enumi}{2}
\item
  If the set \(\mathcal{P}_0\) of functions in \(\mathcal{P}\) that are continuous and finite on X is total in \(\mathcal{C}(\mathrm{X};\mathbf{R})\) , then the relation \(\mu \prec \nu\) is an order relation on \(\mathcal{M}_{+}(\mathrm{X})\) . If a measure \(\nu \in \mathcal{M}_{+}(\mathrm{X})\) is maximal for this order relation, then so is every measure \(\nu'\) such that \(0 \leqslant \nu' \leqslant \nu\) (for the usual order) (argue by contradiction using \(a\) )). If \(\mathcal{P}_0 = \mathcal{P}\) and \(\mathcal{P}\) is total, then every increasing directed set for the order relation \(\mu \prec \nu\) admits a supremum in \(\mathcal{M}_{+}(\mathrm{X})\) for this relation; in particular, \(\mathcal{M}_{+}(\mathrm{X})\) is inductive for this relation.
\item
  Suppose that the functions in \(\mathcal{P}\) are continuous and finite, that \(\mathcal{P}\) is total in \(\mathcal{C}(\mathrm{X};\mathbf{R})\) , and that \(\mathcal{P} + \mathcal{P} \subset \mathcal{P}\) and \(c \cdot \mathcal{P} \subset \mathcal{P}\) for all \(c > 0\) . For every function \(f \in \mathcal{C}(\mathrm{X};\mathbf{R})\) , denote by \(\mathrm{S}(f)\) the lower envelope of the functions \(h \in -\mathcal{P}\) such that \(f \leqslant h\) ; for every \(\varepsilon > 0\) , denote by \(\mathrm{K}_{f,\varepsilon}\) the set of \(x \in \mathrm{X}\) such that \(\mathrm{S}(f)(x) \geqslant f(x) + \varepsilon\) . Show that the following conditions are equivalent:
\end{enumerate}

\(\alpha)\) The measure \(\mu \in \mathcal{M}_{+}(\mathrm{X})\) is maximal for the order relation \(\lambda \prec \lambda'\) .

\(\beta)\) For every function \(f \in \mathcal{C}(\mathrm{X}; \mathbf{R})\) , \(\mu(\mathrm{S}(f)) = \mu(f)\) .

\(\beta^{\prime})\) For every function \(f \in \mathcal{P}\) , \(\mu(\mathrm{S}(f)) = \mu(f)\) .

\(\gamma)\) For every function \(f \in \mathcal{C}(\mathrm{X}; \mathbf{R})\) and every \(\varepsilon > 0\) , \(\mu(\mathrm{K}_{f,\varepsilon}) = 0\) .

\(\gamma^{\prime})\) For every function \(f\in \mathcal{P}\) and every \(\varepsilon >0\) , \(\mu (\mathrm{K}_{f,\varepsilon}) = 0\)

(Apply b) to \(-\mathcal{P}\) .)

\begin{enumerate}
\def\labelenumi{\alph{enumi})}
\setcounter{enumi}{4}
\tightlist
\item
  Under the hypotheses of \(d\) ), show that if \(\mu\) and \(\mu'\) are two positive measures on \(X\) that are maximal for the relation \(\prec\) , then so is \(\mu + \mu'\) . From this, deduce that the set of positive measures that are maximal for this order relation is a lattice for the usual order relation \(\leqslant\) .
\end{enumerate}

¶ 3) Let X be a nonempty compact space, P a set of lower semi-continuous numerical functions on X, taking values in \([-\infty, +\infty]\) ; assume that P contains the finite constants. A point \(x \in X\) is said to be P-extremal if the measure \(\varepsilon_{x}\) is minimal for the pre-order relation \(\mu \prec \nu (*)\) ; the set of P-extremal points is denoted \(\mathrm{Ch}_{\mathcal{P}}(\mathrm{X})\) .

\begin{enumerate}
\def\labelenumi{\alph{enumi})}
\tightlist
\item
  Show that if \(\mathcal{P}'\) is the set of linear combinations \(\sum_{i} c_i h_i\) of functions in \(\mathcal{P}\)
\end{enumerate}

with coefficients \(\geqslant 0\) , the \(\mathcal{P}\) -extremal points are identical to the \(\mathcal{P}'\) -extremal points.

\begin{enumerate}
\def\labelenumi{\alph{enumi})}
\setcounter{enumi}{1}
\tightlist
\item
  For a point \(x \in X\) , show that the following conditions are equivalent:
\end{enumerate}

\(\alpha)\) \(x\) is \(\mathcal{P}\) -extremal.

\(\beta)\) For every function \(f \in \mathcal{C}(\mathrm{X}; \mathbf{R})\) and every \(\varepsilon > 0\) , there exists an \(h \in \mathcal{P}'\) such that \(f \leqslant h\) and \(h(x) \leqslant f(x) + \varepsilon\) .

\(\gamma)\) For every open neighborhood U of \(x\) and every \(\varepsilon > 0\) , there exists a function \(h \geqslant 0\) in \(\mathcal{P}'\) such that \(h(x) \leqslant \varepsilon\) and \(h(y) \geqslant 1\) for all \(y \in X - U\) . (Use Exer. 2 b.)

Show, moreover, that the set of points of \(\mathrm{Ch}_{\mathcal{P}}(\mathrm{X})\) where at least one function in \(\mathcal{P}\) attains its infimum in \(\mathrm{X}\) is dense in \(\mathrm{Ch}_{\mathcal{P}}(\mathrm{X})\) .

\begin{enumerate}
\def\labelenumi{\alph{enumi})}
\setcounter{enumi}{2}
\item
  A subset F of X is said to be P-stable if it is closed and if the relations \(\lambda \prec \mu\) and \(\operatorname{Supp}(\mu) \subset F\) imply \(\operatorname{Supp}(\lambda) \subset F\) . Show that if \(u \in P\) and if \(F'\) is the set of points of F where u attains its infimum in F, then the set \(F'\) is P-stable.
\item
  Assume in addition that \(\mathcal{P}\) separates the points of X. Show that for every function \(h \in \mathcal{P}\) , the set \(S_h\) of points of X where \(h\) attains its infimum in X intersects \(\mathrm{Ch}_{\mathcal{P}}(X)\) . (Consider a minimal \(\mathcal{P}\) -stable subset contained in \(S_h\) and use \(c\) ) to show that such a subset reduces to a single point.)
\item
  With the same hypotheses as in \(d\) ), show that for a closed subset \(\mathbf{F}\) of \(\mathbf{X}\) to contain \(\mathrm{Ch}_{\mathcal{P}}(\mathbf{X})\) , it is necessary and sufficient that for every function \(h \in \mathcal{P}'\) , \(\mathbf{F}\) intersect the set \(\mathbf{S}_h\) of points where \(h\) attains its infimum in \(\mathbf{X}\) (make use of \(b\) ). From this, deduce that for \(a \in \mathbf{X}\) to be in the closure of \(\mathrm{Ch}_{\mathcal{P}}(\mathbf{X})\) , it is necessary and sufficient that for every open neighborhood \(\mathbf{U}\) of \(a\) , there exist \(h \in \mathcal{P}'\) and \(b \in \mathbf{U}\) such that \(h(b) < h(x)\) for all \(x \in \mathbf{X} - \mathbf{U}\) (in other words, \(\mathbf{S}_h \subset \mathbf{U}\) ).
\item
  Under the hypotheses of \(d\) ), a subset A of X is said to be a \(\mathcal{P}\) -edge if, for every function \(h \in \mathcal{P}'\) , \(\mathrm{S}_h\) intersects A. Take for X the union, in \(\mathbf{R}^2\) , of the two circles with respective centers \((-1,0)\) and \((1,0)\) and with radius 1, and for \(\mathcal{P}\) the set of restrictions to X of the affine linear functions on \(\mathbf{R}^2\) . Show that \(\operatorname{Ch}_{\mathcal{P}}(\mathrm{X})\) consists of the points \((\xi, \eta)\) of X such that \(|\xi| \geqslant 1\) . If \(a\) is the \(\mathcal{P}\) -extremal point \((1,1)\) , show that there is a neighborhood U of \(a\) such that there is no function \(h \in \mathcal{P}'\) for which \(h \geqslant 0\) , \(h(a) = 0\) and \(h(y) \geqslant 1\) on \(\mathrm{X} - \mathrm{U}\) . If A and \(\mathrm{A}'\) are the complements in \(\operatorname{Ch}_{\mathcal{P}}(\mathrm{X})\) of the points \((1,1)\) and \((-1,1)\) , respectively, then A and \(\mathrm{A}'\) are \(\mathcal{P}\) -edges but the same is not true of \(\mathrm{A} \cap \mathrm{A}'\) ; thus there is no smallest \(\mathcal{P}\) -edge.
\end{enumerate}

\(g)\) Let \(\mathbf{Y}\) be a second compact space, \(\mathcal{Q}\) a set of lower semi-continuous mappings of \(\mathbf{Y}\) into \([-\infty, +\infty]\) , containing the finite constants. Let \(\mathcal{H}\) be the set of functions \(h = f \otimes g\) with \(f \in \mathcal{P}\) and \(g \in \mathcal{Q}\) ; show that \(\mathrm{Ch}_{\mathcal{H}}(\mathrm{X} \times \mathrm{Y}) = \mathrm{Ch}_{\mathcal{P}}(\mathrm{X}) \times \mathrm{Ch}_{\mathcal{Q}}(\mathrm{Y})\) .

¶ 4) Let E be a Hausdorff locally convex space, X a compact convex set in E, P the set of finite numerical functions continuous and convex in X; \(\mu \prec \nu\) denotes the pre-order relation defined on \(\mathcal{M}_{+}(X)\) by P (Exer. 2).

\begin{enumerate}
\def\labelenumi{\alph{enumi})}
\item
  Show that \(\mu \prec \nu\) is an order relation on \(\mathcal{M}_{+}(\mathrm{X})\) (use Exer. 29 of TVS, II, §5) and that, for every \(x \in \mathrm{X}\) , the relation \(\varepsilon_x \prec \mu\) is equivalent to saying that \(\mu\) has mass 1 and \(x\) is the barycenter of \(\mu\) . For every maximal measure \(\mu\) (for the preceding order relation) of mass 1, there therefore exists one and only one \(x \in \mathrm{X}\) such that \(\varepsilon_x \prec \mu\) . Conversely, every \(z \in \mathrm{X}\) is the barycenter of at least one maximal measure. For \(x \in \mathrm{X}\) to be an extremal point of \(\mathrm{X}\) , it is necessary and sufficient that \(\varepsilon_x\) be a maximal measure.
\item
  Let \(\mu\) be a maximal measure on \(\mathrm{X}\) , \(f\) a function in \(\mathcal{C}_{+}(\mathrm{X})\) . Show that for every \(\varepsilon > 0\) , there exists a continuous convex function \(g\) on \(\mathrm{X}\) such that \(0 \leqslant g \leqslant f\) and \(\mu(g) \geqslant \mu(f) - \varepsilon\) (use Exers. 3 b) and 2 d)). From this, deduce that the support of \(\mu\) is contained in the closure of the set of extremal points of \(\mathrm{X}\) .
\item
  Let \(\mu\) be a maximal measure on \(X\) , \((f_n)_{n \geqslant 1}\) a decreasing sequence of functions in \(\mathcal{C}_+(X)\) . Show that if, for every extremal point \(x \in X\) , the sequence \((f_n(x))\) tends to 0, then \(\lim_{n \to \infty} \mu(f_n) = 0\) (use \(b\) ) to construct a decreasing sequence \((g_n)\) of continuous convex functions such that \(0 \leqslant g_n \leqslant f_n\) and \(\mu(g_n) \geqslant \mu(f_n) - \varepsilon\) for all \(n\) , and show that the sequence \((g_n(y))\) tends to 0 for every \(y \in X\) ).
\item
  Deduce from c) that if \(A \subset X\) contains no extremal point and is the union of a sequence \((\mathrm{K}_{n})\) of compact subsets of X, each of which is a countable intersection of open sets in X, then \(\mu(\mathrm{A}) = 0\) for every maximal measure \(\mu\) .
\item
  Let \(\mathcal{A}\) be the vector space of continuous functions on X that are the restrictions to X of continuous affine linear functions on E. The relation \(\lambda \prec \mu\) implies that \(\lambda(h) = \mu(h)\) for every function \(h \in \mathcal{A}\) ; for every function \(f \in \mathcal{C}(\mathrm{X};\mathbf{R})\) , the lower envelope \(\mathrm{S}(f)\) of the functions \(h \in \mathcal{A}\) such that \(f \leqslant h\) is also the lower envelope of the functions \(g \in -\mathcal{P}\) such that \(f \leqslant g\) . For every \(x \in X\) , the relation \(\varepsilon_x \prec \mu\) is equivalent to the relation \(h(x) = \langle h, \mu \rangle\) for all \(h \in \mathcal{A}\) ; the extremal points of X are identical to the \(\mathcal{A}\) -extremal points. If \(f \in \mathcal{P}\) then, for every \(x \in X\) ,
\end{enumerate}

\[
\mathrm{S} (f) (x) = \sup _ {\varepsilon_ {x} \prec \mu} \int f (y) d \mu (y)\tag{*}
\]

(use Exer. 2 b)). This formula is no longer necessarily valid if it is only assumed that f is convex and upper semi-continuous on X (with the notations of Exer. 1 b), consider the function f equal to 1 on the image of I in K, and to 0 elsewhere).

\begin{enumerate}
\def\labelenumi{\alph{enumi})}
\setcounter{enumi}{5}
\tightlist
\item
  Show that the formula (*) is again valid when in the second member one limits oneself to the discrete measures \(\mu\) having x as barycenter.
\end{enumerate}

¶ 5) The notations being those of Exer. 4 e), let \(\mathcal{A}'\) denote the dual of the normed space \(\mathcal{A}\) , equipped with the order structure deduced from that of \(\mathcal{A}\) (Ch. II, §2). On the other hand, denote by C the convex cone with vertex 0 generated by \(X \times \{1\}\) in the space \(E \times R\) .

\begin{enumerate}
\def\labelenumi{\alph{enumi})}
\tightlist
\item
  Show that the following properties are equivalent:
\end{enumerate}

\begin{enumerate}
\def\labelenumi{(\roman{enumi})}
\item
  Every point \(x \in X\) is the barycenter of a unique maximal measure \(\beta_x\) of mass 1 on \(X\) .
\item
  The vector space F generated by C is a lattice for the order structure having C as the set of elements \(\geqslant0\) (in other words, X is a simplex (TVS, II, §2, Exer. 41)).
\item
  For every function \(f \in \mathcal{P}\) , \(S(f)\) is both concave and convex.
\item
  If \(\mu\) and \(\nu\) are two maximal measures such that \(\mu(h) = \nu(h)\) for all \(h \in \mathcal{A}\) , then \(\mu = \nu\) .
\item
  The vector space \(\mathcal{A}'\) is a Riesz space.
\item
  If \(f\) and \(g\) are two functions in \(\mathcal{P}\) , then \(S(f + g) = S(f) + S(g)\) .
\end{enumerate}

(First prove that (i) implies the following property:

\begin{enumerate}
\def\labelenumi{(\roman{enumi})}
\setcounter{enumi}{6}
\tightlist
\item
  The mapping \(x \mapsto \beta_{x}\) may be extended in only one way to a bijective linear mapping of F onto the subspace of \(\mathcal{M}(X)\) generated by the maximal measures, and this mapping transforms C into the cone of maximal measures.
\end{enumerate}

Then deduce (ii) from (vii), using Exer. 2 e). To deduce (iii) from (ii), use the decomposition lemma in the Riesz space F and Exer. 4 f). To deduce (iv) from (iii), use the fact that if \(g \in \mathcal{C}(X; \mathbf{R})\) is both concave and convex, then the set of \(h \in A\) such that \(h(x) > g(x)\) for all \(x \in X\) is a decreasing directed set (TVS, II, §5, Prop. 6) and has g for its lower envelope; then apply Exer. 4 e). To deduce (v) from (iv), observe that \(A'\) may be identified with the subspace of \(\mathcal{M}(X)\) generated by the maximal measures, with the help of Exer. 4 e). To deduce (vi) from (v), apply the decomposition lemma in the Riesz space \(A'\) and Exer. 4 f). Finally, to deduce (i) from (vi), consider the mapping \(f \mapsto S(f)(x)\) for \(f \in P\) and \(x \in X\) , and use Exer. 4 e).

\begin{enumerate}
\def\labelenumi{\alph{enumi})}
\setcounter{enumi}{1}
\item
  When the conditions of a) are satisfied, show that for every function \(f \in \mathcal{C}(X; \mathbf{R})\) , the mapping \(x \mapsto \beta_x(f)\) is in the closure, for the topology of uniform convergence in X, of the set of bounded functions on X that are the difference of two upper semi-continuous functions (make use of the fact that P is total in \(\mathcal{C}(X; \mathbf{R})\) ); consequently, this mapping is measurable for every measure on X. *For every positive measure \(\mu\) on X, the unique maximal measure \(\nu \succ \mu\) is given by \(\nu = \int \beta_x d\mu(x)\) (cf.~Ch. V).*
\item
  When the conditions of \(a\) ) are satisfied and X is in addition metrizable, show that for every \(x \in X\) , \(\beta_x\) is the only positive measure \(\mu\) of mass 1 and barycenter \(x\) such that \(\mu(X - L) = 0\) , where \(L\) is the set \(\mathrm{Ch}_{\mathcal{A}}(X)\) of extremal points (argue as in Th. 1).
\item
  When the conditions of a) are satisfied, show that the following properties are equivalent:
\end{enumerate}

\(\alpha)\) The set \(\mathrm{L} = \mathrm{Ch}_{\mathcal{A}}(\mathrm{X})\) is closed, in other words \(\check{\mathrm{S}}_{\mathcal{A}}(\mathrm{X}) = \mathrm{Ch}_{\mathcal{A}}(\mathrm{X})\) .

\(\beta)\) For every function \(f \in \mathcal{P}\) , the restriction of \(\mathrm{S}(f)\) to \(\overline{\mathrm{L}}\) is continuous.

\(\gamma)\) For every function \(f \in \mathcal{P}\) , the function \(S(f)\) is continuous on \(X\) .

\(\delta)\) The mapping \(x \mapsto \beta_x\) of X into \(\mathcal{M}_+(X)\) (equipped with the vague topology) is continuous.

(To show that \(\alpha\) ) implies \(\beta\) ), observe that \(\beta_x = \varepsilon_x\) in L. To prove that \(\beta\) implies \(\gamma\) ), use the fact that S(f) is the lower envelope of a decreasing directed set of functions in \(\mathcal{A}\) (TVS, II, §5, Prop. 6), Dini's theorem and Exer. 3 \(d\) ) applied to the functions in \(\mathcal{A}\) . To prove that \(\gamma\) implies \(\delta\) ), use Exer. 4 \(e\) ). Finally, to see that \(\delta\) implies \(\alpha\) ), use Exer. 4 \(a\) .)

¶ 6) Let X be a nonempty compact space, H a linear subspace of \(\mathcal{C}(X;R)\) containing the constants and separating the points of X. The relation \(\lambda \prec \mu\) defined by H (Exer. 2) is then an equivalence relation, and the points H-extremal in the sense of Exer. 3 are identical to the points H-extremal in the sense of No.~3.

Let \(\mathbf{F}\) be a closed subset of \(\mathbf{X}\) containing \(\operatorname{Ch}_{\mathcal{H}}(\mathbf{X})\) ; for every \(x \in \mathbf{X}\) , denote by \(\mathcal{M}_x^{\mathrm{F}}\) the set of positive measures \(\mu\) on \(\mathbf{X}\) , of total mass 1, such that \(\varepsilon_x \prec \mu\) and \(\operatorname{Supp}(\mu) \subset \mathbf{F}\) ; in order that \(\varepsilon_x \in \mathcal{M}_x^{\mathrm{F}}\) , it is necessary and sufficient that \(x \in \mathbf{F}\) ; the relation \(x \in \operatorname{Ch}_{\mathcal{H}}(\mathbf{X})\) is equivalent to \(\mathcal{M}_x^{\mathrm{F}} = \{\varepsilon_x\}\) . For every bounded numerical function \(f\) defined on \(\mathbf{F}\) and for every \(x \in X\) , set

\[
\overline {{\mathrm{H}}} _ {x} ^ {\mathrm{F}} (f) = \inf _ {f \leqslant h | \mathrm{F}, h \in \mathscr {H}} h (x), \quad \underline {{\mathrm{H}}} _ {x} ^ {\mathrm{F}} (f) = - \overline {{\mathrm{H}}} _ {x} ^ {\mathrm{F}} (- f) = \sup _ {f \geqslant h | \mathrm{F}, h \in \mathscr {H}} h (x),
\]

which are finite numbers.

\begin{enumerate}
\def\labelenumi{\alph{enumi})}
\item
  Show that the mapping \(f \mapsto \overline{\mathrm{H}}_x^{\mathrm{F}}(f)\) of \(\mathcal{C}(\mathrm{F};\mathbf{R})\) into \(\mathbf{R}\) is increasing, positively homogeneous and convex. If \(h \in \mathcal{H}\) , then \(\overline{\mathrm{H}}_x^{\mathrm{F}}(h) = h(x)\) .
\item
  For every \(x \in X\) , every measure \(\mu \in \mathcal{M}_x^{\mathrm{F}}\) and every function \(f \in \mathcal{C}(\mathrm{F};\mathbf{R})\) , \(\underline{\mathrm{H}}_x^{\mathrm{F}}(f) \leqslant \int f d\mu \leqslant \overline{\mathrm{H}}_x^{\mathrm{F}}(f)\) . Conversely if, for a function \(f_0 \in \mathcal{C}(\mathrm{F};\mathbf{R})\) , \(\gamma\) is a real number such that \(\underline{\mathrm{H}}_x^{\mathrm{F}}(f_0) \leqslant \gamma \leqslant \overline{\mathrm{H}}_x^{\mathrm{F}}(f_0)\) , then there exists a measure \(\mu \in \mathcal{M}_x^{\mathrm{F}}\) such that \(\gamma = \int f_0 d\mu\) .
\item
  Show that for a function \(g \in \mathcal{C}(\mathrm{X};\mathbf{R})\) , the following conditions are equivalent:
\end{enumerate}

\(\alpha)\) For every \(x \in X\) and every measure \(\mu \in \mathcal{M}_x^{\mathrm{F}}\) , one has \(g(x) = \int g d\mu\) .

\(\beta)\) For every \(x \in \mathrm{X}\) , \(\underline{\mathrm{H}}_x^{\mathrm{F}}(g|\mathrm{F}) = \overline{\mathrm{H}}_x^{\mathrm{F}}(g|\mathrm{F}) = g(x)\) .

\(\gamma)\) For every \(\varepsilon > 0\) there exist two finite sequences \((h_i'), (h_j'')\) of functions in \(\mathcal{H}\) such that, setting \(h' = \sup(h_i')\) , \(h'' = \inf(h_j'')\) , one has \(h' \leqslant g \leqslant h''\) and \(h'' - h' \leqslant \varepsilon\) .

(To see that \(\beta\) ) implies \(\gamma\) ), argue as in GT, X, §4, No.~1, Prop. 2.)

When a function \(g \in \mathcal{C}(\mathrm{X};\mathbf{R})\) has the preceding equivalent properties, it is said to be \(\mathcal{H}\) -harmonic; the set \(\mathcal{H}^c\) of \(\mathcal{H}\) -harmonic functions is a closed linear subspace of \(\mathcal{C}(\mathrm{X};\mathbf{R})\) containing \(\mathcal{H}\) ; it is independent of the closed set \(\mathrm{F} \supset \mathrm{Ch}_{\mathcal{H}}(\mathrm{X})\) under consideration. One has \((\mathcal{H}^c)^c = \mathcal{H}^c\) , and the equivalence relation in \(\mathcal{M}_+(\mathrm{X})\) defined by \(\mathcal{H}^c\) is identical to the relation \(\lambda \prec \mu\) defined by \(\mathcal{H}\) ; consequently \(\mathrm{Ch}_{\mathcal{H}^c}(\mathrm{X}) = \mathrm{Ch}_{\mathcal{H}}(\mathrm{X})\) . The mapping \(g \mapsto g|\mathrm{F}\) is a strictly increasing isometry of \(\mathcal{H}^c\) onto its image in \(\mathcal{C}(\mathrm{F};\mathbf{R})\) (which is therefore a closed subspace of \(\mathcal{C}(\mathrm{F};\mathbf{R})\) ).

\begin{enumerate}
\def\labelenumi{\alph{enumi})}
\setcounter{enumi}{3}
\item
  Take for X the interval \([0,1]\) of R, and for H the vector space of restrictions to X of the polynomials of second degree over R. Show that \(H^{c}\) is distinct from the closure \(\overline{H}\) of H in \(\mathcal{C}(X;R)\) .
\item
  Show that in order that \(\mathrm{Ch}_{\mathcal{H}}(\mathrm{X}) = \mathrm{X}\) , it is necessary and sufficient that \(\mathcal{H}^c = \mathcal{C}(\mathrm{X};\mathbf{R})\) (to prove that the condition is necessary, use \(b\) )).
\item
  Show that if \(\mathcal{H}\) is lattice-ordered (in other words, is a Riesz space), then \(\mathcal{H}^c = \overline{\mathcal{H}}\) . Give an example where \(\mathcal{H}\) is not lattice-ordered and \(\mathcal{H}^c = \overline{\mathcal{H}}\) .
\item
  Let \(\mathcal{E}_{\mathcal{H}}\) be the smallest closed subset of \(\mathcal{C}(\mathrm{X};\mathbf{R})\) containing \(\mathcal{H}\) such that the lower envelope \(\inf (u,v)\) of any two functions of \(\mathcal{E}_{\mathcal{H}}\) belongs to \(\mathcal{E}_{\mathcal{H}}\) . Show that the following properties are equivalent:
\end{enumerate}

\(\alpha)\) \(f\in \mathcal{E}_{\mathcal{H}}\)

\(\beta)\) for every \(x \in X\) and every measure \(\mu \in \mathcal{M}_x^X\) , \(\int f d\mu \leqslant f(x)\) ;

\(\gamma)\ \overline{\mathrm{H}}_{x}^{\mathrm{X}}(f) = f(x)\ \text{for every } x \in \mathrm{X};\)

\(\delta)\) for every \(\varepsilon > 0\) , there exists a finite sequence \((h_i)\) of functions in \(\mathcal{H}\) such that \(f \leqslant \inf(h_i) \leqslant f + \varepsilon\) .

(To see that \(\beta\) ) implies \(\gamma\) , make use of \(b\) .)

Deduce from this that \(\mathcal{E}_{\mathcal{H}}\) is a pointed convex cone, and that \(\mathcal{E}_{\mathcal{H}} \cap (-\mathcal{E}_{\mathcal{H}}) = \mathcal{H}^c\) . Show that every function in \(\mathcal{E}_{\mathcal{H}}\) attains its infimum in \(X\) at at least one point of \(\mathrm{Ch}_{\mathcal{H}}(X)\) .

¶ 7) Let Y be a nonempty compact space, R a linear subspace of \(\mathcal{C}(\mathrm{Y};\mathbf{R})\) that contains the constants, separates the points of Y and is lattice-ordered (in other words, is a Riesz space).

\begin{enumerate}
\def\labelenumi{\alph{enumi})}
\tightlist
\item
  Let \(\mathcal{N}\) be a maximal isolated subspace of \(\mathcal{R}\) (Ch. II, §1, Exer. 4); show that there exists one and only one point \(y_0\) such that the positive linear form \(\varphi : f \mapsto f(y_0)\)
\end{enumerate}

on \(\mathcal{R}\) is latticial (Ch. II, §2, Exer. 5 b)) and \(\mathcal{N} = \overline{\varphi}(0)\) . (To see that there exists at least one point of Y where all of the functions \(f \geqslant 0\) belonging to \(\mathcal{N}\) are zero, argue by contradiction, by showing that in the contrary case, taking into account the compactness of Y and the definition of the order relation in R, one would have N = R. To show that the set \(Z(N)\) of points of Y where all of the functions in N are zero reduces to a single point, make use of the fact that N is a hyperplane in R (Ch. II, §2, Exer. 5) and the fact that R separates the points of Y and contains the constants.)

\begin{enumerate}
\def\labelenumi{\alph{enumi})}
\setcounter{enumi}{1}
\item
  Under the hypotheses of \(a\) ), show that \(y_0 \in \check{\mathrm{S}}_{\mathcal{R}}(\mathrm{Y})\) . (Setting \(\mathrm{S} = \check{\mathrm{S}}_{\mathcal{R}}(\mathrm{Y})\) , note that the subspace \(\mathcal{R}'\) of \(\mathcal{C}(\mathrm{S}; \mathbf{R})\) formed by the restrictions to \(\mathrm{S}\) of the functions in \(\mathcal{R}\) is canonically isomorphic to \(\mathcal{R}\) as an ordered vector space (Exer. 6 c)) and that, by means of the isomorphism inverse to \(f \mapsto f|\mathrm{S}\) , the positive linear form \(f \mapsto f(y_0)\) yields a latticial positive linear form on \(\mathcal{R}'\) ; then apply \(a\) ) to \(\mathcal{R}'\) .)
\item
  For every isolated subspace \(\mathcal{N}_0 \neq \mathcal{R}\) , show that there exists a maximal isolated subspace \(\mathcal{N} \supset \mathcal{N}_0\) . (Use the fact that if an isolated subspace contains a constant function \(\neq 0\) , then it is equal to \(\mathcal{R}\) .)
\item
  Let \(\mathbf{Z}\) be the union of the sets \(\mathbf{Z}(\mathcal{N})\) (each reduced to a point) as \(\mathcal{N}\) runs over the set of maximal isolated subspaces of \(\mathcal{R}\) ; show that \(\mathbf{Z} = \check{\mathbf{S}}_{\mathcal{R}}(\mathbf{Y})\) . (First prove that \(\mathbf{Z}\) is closed, by noting that if \(y \notin \mathbb{Z}\) then the linear form \(f \mapsto f(y)\) on \(\mathcal{R}\) is not latticial, consequently (using Ch. II, §2, Exer. 5 b)) neither is the linear form \(f \mapsto f(y')\) for all points \(y'\) sufficiently near \(y\) . Next, use Prop. 7 of No.~3, by showing that every function \(f_0 \in \mathcal{R}\) attains its infimum \(\alpha\) in \(\mathbf{Y}\) at at least one point of \(\mathbf{Z}\) ; for this, setting \(\mathrm{M} = f_0^{-1}(\alpha)\) , one considers the subspace \(\mathcal{N}_0\) of functions \(f \in \mathcal{R}\) of the form \(f_1 - f_2\) .
\end{enumerate}

\(\mathrm{M} = \overline{f}_0(\alpha)\) , one considers the subspace \(\mathcal{N}_0\) of functions \(f \in \mathcal{R}\) of the form \(f_1 - f_2\) with \(f_1 \geqslant 0\) , \(f_2 \geqslant 0\) , \(f_1\) and \(f_2\) zero on \(\mathrm{M}\) ; prove that \(\mathcal{N}_0\) is isolated and use \(c\) .)

¶ 8) The general notations and hypotheses being those of Exer. 6, set \(S = \check{S}_{\mathcal{H}}(X)\) . A function \(f \in \mathcal{C}(S; \mathbf{R})\) is said to be \(\mathcal{H}\) -resolvent if there exists a function \(g \in \mathcal{H}^c\) such that \(f = g|S\) (this function is then unique).

\begin{enumerate}
\def\labelenumi{\alph{enumi})}
\tightlist
\item
  Show that the following conditions on a function \(f \in \mathcal{C}(\mathbf{S};\mathbf{R})\) are equivalent:
\end{enumerate}

\begin{enumerate}
\def\labelenumi{(\roman{enumi})}
\item
  \(f\) is \(\mathcal{H}\) -resolvent.
\item
  For every \(x \in X\) , \(\underline{\mathrm{H}}_x^{\mathrm{S}}(f) = \overline{\mathrm{H}}_x^{\mathrm{S}}(f)\) .
\item
  For every \(x \in X\) and every pair of measures \(\mu_1, \mu_2\) in \(\mathcal{M}_x^S\) , one has \(\int f d\mu_1 = \int f d\mu_2\) .
\end{enumerate}

(Make use of Exer. 6 b) and c).)

\begin{enumerate}
\def\labelenumi{\alph{enumi})}
\setcounter{enumi}{1}
\tightlist
\item
  Show that the following properties are equivalent:
\end{enumerate}

\begin{enumerate}
\def\labelenumi{(\roman{enumi})}
\item
  Every function \(f \in \mathcal{C}(\mathrm{S};\mathbf{R})\) is \(\mathcal{H}\) -resolvent.
\item
  For every \(x \in X\) and every function \(f \in \mathcal{C}(\mathrm{S};\mathbf{R})\) , \(\underline{\mathrm{H}}_x^{\mathrm{S}}(f) = \overline{\mathrm{H}}_x^{\mathrm{S}}(f)\) .
\item
  The set \(\mathcal{M}_x^{\mathrm{S}}\) reduces to a single element.
\item
  The set \(\mathcal{H}^{\tilde{c}}\) is lattice-ordered.
\item
  For every function \(u \in \mathcal{E}_{\mathcal{H}}\) (Exer. 6 g)), there exists a greatest lower bound \(h_u\) of \(u\) in \(\mathcal{H}^c\) .
\end{enumerate}

(To show that (iv) implies (i), apply Exer. 7 d) to \(\mathcal{H}^c\) , proving that for every point \(s \in S\) , the linear form \(h \mapsto h(s)\) is latticial in \(\mathcal{H}^c\) ; then use Stone's theorem and Exer. 6 c). To see that (i) implies (v), note that \(\mathcal{H}^c \subset \mathcal{E}_{\mathcal{H}}\) and consider the unique function \(h_u \in \mathcal{H}^c\) such that \(h_u|S = u|S\) . To see that (v) implies (iv), note that if \(h \in \mathcal{H}^c\) then \(-h^{-} = \inf(h,0) \in \mathcal{E}_{\mathcal{H}}\) , and apply (v) to \(u = -h^{-}\) .)

\begin{enumerate}
\def\labelenumi{\alph{enumi})}
\setcounter{enumi}{2}
\item
  If the conditions of \(b\) ) are satisfied and if \(\gamma_x\) is the unique element of \(\mathcal{M}_x^S\) , show that the mapping \(x \mapsto \gamma_x\) of \(X\) into \(\mathcal{M}_+(X)\) (equipped with the vague topology) is continuous; from this, deduce that one then has \(\mathrm{Ch}_{\mathcal{H}}(X) = \check{\mathrm{S}}_{\mathcal{H}}(X)\) (note that \(\gamma_x = \varepsilon_x\) in \(\mathrm{Ch}_{\mathcal{H}}(X)\) ).
\item
  In order that every function \(f \in \mathcal{C}(\mathrm{S};\mathbf{R})\) be the restriction to \(S\) of a function \(h \in \mathcal{H}\) , it is necessary and sufficient that \(\mathcal{H}\) be closed in \(\mathcal{C}(\mathrm{X};\mathbf{R})\) and be lattice-ordered.
\end{enumerate}

\begin{enumerate}
\def\labelenumi{\arabic{enumi})}
\setcounter{enumi}{8}
\tightlist
\item
  Let Y be the topological space whose underlying set is the product \(I \times \{-1,0,1\}\) in \(R^{2}\) , where \(I = [0,1]\) ; for every \(a \in I\) and every \(\varepsilon > 0\) , let \(U_{a,\varepsilon}\) be the set of \((x,y) \in Y\) such that \(|x - a| \leqslant \varepsilon\) and \((x,y)\) is distinct from \((a,1)\) and \((a,-1)\) ; the sets \(\{(x,y)\}\) for \(y \neq 0\) and \(x \in I\) together with the sets \(U_{a,\varepsilon}\) form a base for a topology on Y for which Y is a non-metrizable compact space (cf.~GT, IX, §2, Exer. 13 d)). Denote by X the topological sum space of Y and a set reduced to a point \(\omega\) . Let \(\mathcal{H}\) be the set of continuous finite numerical functions \(h\) on X such that
\end{enumerate}

\[
h (x, 0) = \frac {1}{2} \bigl (h (x, 1) + h (x, - 1) \bigr)
\]

for all \(x \in I\) , and \(h(\omega) = \int_0^1 h(x, 0) dx\) . Show that \(\operatorname{Ch}_{\mathcal{H}}(X)\) consists of the points \((x, y)\) such that \(y \neq 0\) and \(x \in I\) , but that there exists no positive measure \(\mu\) of mass 1 on X such that \(h(\omega) = \mu(h)\) for all \(h \in \mathcal{H}\) , and \(\mu(X - \operatorname{Ch}_{\mathcal{H}}(X)) = 0\) .

¶ 10) a) Let E be a Hausdorff and complete locally convex space, E' its dual, and E'* ⊃ E the algebraic dual of E'. Let A be a subset of E that is compact for the weakened topology σ(E, E') . Let x be a point of E'* in the closure of the convex envelope C of A for the weak topology σ(E',E'). Show that if \((x'_n)\) is a sequence of points of E' that converges to a' ∈ E' for the weak topology σ(E',E), then \(\lim_{n \to \infty}\langle x, x'_n\rangle = \langle x, a'\rangle\).

(Note that \(x\) is the barycenter of a measure of mass 1 on A, and apply Lebesgue's theorem.)

\begin{enumerate}
\def\labelenumi{\alph{enumi})}
\setcounter{enumi}{1}
\item
  Deduce from a) that if there exists in E a sequence of points that is dense for the original topology, then the restriction to every equicontinuous subset \(H'\) of \(E'\) of the mapping \(x' \mapsto \langle x, x' \rangle\) is continuous for the weak topology \(\sigma(\mathrm{E}', \mathrm{E})\) (note that the topology induced on \(H'\) by \(\sigma(\mathrm{E}', \mathrm{E})\) is metrizable). Deduce that in this case, necessarily \(x \in E\) (TVS, IV, §5, No.~5); in other words, the closed convex envelope in E of a set that is compact for the weakened topology is again compact for that topology.
\item
  Extend the result of \(b\) ) to the case that \(E\) is any quasi-complete Hausdorff locally convex space (Krein's theorem). (First reduce to the case that \(E\) is complete by considering \(\widehat{E}\) ; then note that by virtue of Eberlein's theorem (TVS, IV, §5, No.~3, Th. 1), it suffices to prove that every sequence \((x_n)\) of points of \(C\) has a cluster point in \(E\) for \(\sigma(E, E')\) ; this permits reducing to the case that there exists in \(E\) a sequence dense for the original topology.)
\end{enumerate}

